\section{Conclusion}
\label{sec:conclusion}

We built a single simulation engine in which nine clustering protocols spanning
twenty-five years of design practice execute under provably identical physical
conditions, and evaluated them over 600 runs across two channel configurations
and a further 1{,}350 runs across nine deployment scales. The engine owns energy
accounting, packet sizing, fusion charges, the channel and reading accounting; a
protocol returns only a cluster structure. Fairness is therefore a property of
the construction rather than of the authors' care, and it is verified by static
audit rather than asserted.

Five findings hold up.

Clustering does not extend network lifetime so much as redistribute it. LEACH
delays first node death by a factor of 9.2 over direct transmission while
bringing last death forward by nearly a third, because about a third of the
nodes lie close enough to the sink to transmit on the free-space branch and are
penalized rather than helped by rotation. No single lifetime figure summarizes
this, and the two most commonly reported figures order the protocols in opposite
directions.

The advantage of the soft-computing and learned protocols over LEACH in first
node death is conditional on the channel. Under packet loss the fuzzy system and
the Q-network improve on LEACH by 417 and 426 rounds; with loss removed the same
comparisons give 17 and 22 rounds and lose significance. Direct measurement of
the energy split shows why, and the explanation is not the one the mean
distances suggest. Every protocol's mean head-to-sink distance sits below 120~m,
where packet error is still effectively zero, so the mean cannot be the cause.
The tails differ instead: LEACH puts 29.9\% of its head-rounds beyond 120~m and
spends 6.39\% of its budget on retransmissions, against the Q-network's 5.1\% and
2.31\%. The demonstrable skill of the newer protocols is keeping the far tail of
the sink hop out of the error waterfall, and it pays in proportion to how lossy
that hop is. Their advantage in overall survival and in readings delivered is not
conditional; it holds in both configurations at the corrected significance floor.

The same boundary appears from an independent direction when the deployment
scales. Across three node counts and three field areas, tripling the node count
moves first node death by a factor of 0.88 to 1.26 while tripling the field side
moves it by a factor of 5 to 7, and cells of nearly equal density differ by up to
5.5 times. Distance to the sink is the variable, not density. In a
$50\times50$~m field, where nothing reaches the waterfall, the learned and
soft-computing advantages collapse exactly as they do when packet loss is
switched off, and clustering itself becomes counterproductive: the no-clustering
baseline outlives LEACH in all three of those cells.

Temporal credit assignment is worth measuring separately from representational
capacity. The Q-network and the graph network consume identical features and
differ only in architecture and training signal. The Q-network's first-death
distribution spans 219 rounds; the graph network's spans 388 and is bimodal, and
its busiest node serves as head 13.6 times as often as LEACH's while 13 nodes
never serve at all. A single-round objective cannot represent deferring cost to
extend lifetime even when rotation state is visible in its inputs.

Reactive reporting buys lifetime at a price that must be reported alongside it.
TEEN and APTEEN survive past the round cap in every trial with an area under the
living-node curve more than three times LEACH's, at a data yield of 0.11 to 0.12
against roughly 0.9 for the proactive protocols.

The methodological point generalizes beyond these nine protocols. Running the
entire study at both endpoints of an unspecified parameter, and committing in
advance to report only what survives both, retracted a comparison that either
endpoint alone would have supported. Simulation studies routinely fix such
parameters silently; where the parameter sits near a nonlinearity, as transmit
power does relative to the packet-error waterfall, the choice can decide the
ranking. The cost of bracketing it was a doubled runtime, which at this scale is
eighty minutes.

The largest confound remaining is the accounting subsidy granted to centralized
protocols, worth nine to ten percentage points of the energy budget. An ablation
that charges the uplink honestly would convert it from a caveat into a
measurement, and is the most valuable extension of this work. Sweeping transmit
power across several intermediate values would trace the curve that the two
channel endpoints only bracket.

All configuration constants, per-round outputs, verification artifacts and the
packet-error calibration curve are released with the code.
